\chapter{Introducción}
\section{Introducción General}

En la zona central de Chile, el 95 \% de la sismicidad registrada está asociada al proceso de subducción de la placa de Nazca bajo la placa Sudamericana \shortcite{Ammirati2019}. Esta región sismogénica ha producido eventos sísmicos, como el terremoto de Maule en 2010, con una magnitud de \(\Mw{}\,8.8\) \shortcite{Ruiz2018}. La mayoría de estos eventos corresponden a sismos interplaca \shortcite{Ruiz2018}, aunque también se han registrado sismos intraplaca, como el terremoto de Chillán en 1939, con magnitud de \(\Mw{}\,7.8\) \shortcite{Leyton2009}, y el terremoto de Punitaqui en 1997, con una magnitud de \(\Mw{}\,7.1\) \shortcite{Pardo1997}. Asimismo, se han reportado sismos corticales, como el terremoto de Las Melosas en 1958, con una magnitud de \(M_s\,6.9\), y el de Curicó en 2004, con una magnitud de \(\Mw{}\,6.4\) \shortcite{Campos2004}.

En el piedemonte de la cordillera de los Andes, en Santiago, se encuentra la Falla San Ramón (FSR), un sistema de falla cortical inversa que se extiende aproximadamente 50 km en dirección norte-sur, donde el plano de falla alcanza la superficie \shortcite{Ammirati2019}. Este tipo de estructuras representa una fuente de actividad sísmica que puede afectar directamente a las áreas urbanas adyacentes. Dada su proximidad a la ciudad, la FSR tiene un alto potencial de impacto.

A nivel mundial, los sismos corticales superficiales han causado impactos significativos en áreas densamente pobladas. Un ejemplo notable es el terremoto de Chi-Chi en 1999, de magnitud \(\Mw{}\,7.6\), que provocó 2470 muertes y destruyó más de 100,000 estructuras \shortcite{Shin2001}. Otros casos incluyen el terremoto de Kobe en 1995, de magnitud \(\Mw{}\,6.9\) \shortcite{Kanaori1996}, y el terremoto de Christchurch en 2010, de magnitud \(\Mw{}\,7.1\) \shortcite{Quigley2016}. Estos eventos destacan los posibles efectos devastadores de sismos similares en áreas urbanas, como Santiago.

Con una población de 7,112,808 habitantes \shortcite{ine_censo2017} y un PIB de 144,356 millones de USD en 2023 \shortcite{BancoCentralChile2023}, Santiago se encuentra en una región donde confluyen diversas fuentes sismogénicas. La interacción entre la densidad poblacional, la actividad sísmica histórica y las características urbanas subrayan la importancia de comprender mejor estas amenazas para mitigar riesgos.

Dado que los escarpes de la FSR presentan diversos grados de erosión que han alterado su morfología, se ha establecido, para la modificación del Plan Regulador Metropolitano de Santiago (PRMS), un rango teórico en el cual podría producirse la ruptura superficial de la falla respecto a su traza oficial. Este rango abarca 300 metros: 100 metros hacia el oeste y 200 metros hacia el este \shortcite{seremi2012riesgo}. Como resultado, surge la necesidad de realizar un estudio de riesgo detallado para las estructuras ubicadas dentro de esta zona buffer.

En los últimos años, se han realizado estudios de amenaza asociados a la FSR. Estos análisis han estimado que las magnitudes máximas que podría generar oscilan entre \(\Mw{}\,6.9\) y \(\Mw{}\,7.5\), mientras que las aceleraciones cercanas al escarpe de la falla se sitúan en un rango de 0.7 a 0.8 g \shortcite{Armijo2010,Estay2016,Ammirati2019}. Como referencia, durante el terremoto de Maule en 2010, de magnitud \(\Mw{}\,8.8\), se registró un PGA de \textasciitilde 0.3 g en Santiago Oriente \shortcite{Boroschek2012}.

\section{Hipótesis}

La Falla San Ramón, como fuente sísmica de tipo cortical, podría generar sismos de intensidades suficientes para causar daños materiales y pérdidas económicas en las áreas cercanas a su ubicación.

\section{Objetivos}
\subsection{Objetivo general}

El objetivo general de este trabajo es analizar los posibles escenarios sísmicos asociados a la Falla San Ramón, evaluando las máximas aceleraciones espectrales que podrían generarse durante un eventual terremoto. Asimismo, se desarrollará un modelo para estimar las pérdidas en las zonas cercanas a la FSR.

\subsection{Objetivos específicos}

Los objetivos específicos de esta investigación son:

\begin{enumerate}
    \item \textbf{Recopilar y analizar modelos de movimiento de suelo}, aplicables a la realidad sísmica y geológica de la región estudiada, interpretando su relevancia para el análisis.
    \item \textbf{Desarrollar un modelo tridimensional de la Falla San Ramón (FSR)} en profundidad, basado en registros de sismicidad medidos en la zona de estudio.
    \item \textbf{Elaborar y simular un modelo de amenaza sísmica} utilizando el software OpenQuake.
    \item \textbf{Recopilar y analizar modelos de fragilidad estructural}, relevantes para el contexto chileno, teniendo en cuenta los sistemas estructurales y los materiales predominantes en la ciudad de Santiago, y desarrollar a partir de ello un modelo de vulnerabilidad.
    \item \textbf{Realizar un análisis de riesgo sísmico}, integrando los modelos de amenaza y vulnerabilidad en el software OpenQuake para simular escenarios de riesgo.
\end{enumerate}

\section{Estructura del trabajo}

El trabajo está estructurado en seis capítulos, organizados de la siguiente manera:

\begin{itemize}
    \item \textbf{Capítulo 1: Introducción.} Este capítulo presenta el contexto general, la motivación y los propósitos principales del trabajo.
    \item \textbf{Capítulo 2: Marco Teórico.} Se explican los conceptos teóricos y los supuestos utilizados en el desarrollo de los modelos.
    \item \textbf{Capítulo 3: Contexto Geológico y Sísmico de la Cuenca de Santiago.} Aquí se detallan las características geológicas relevantes que enmarcan el análisis.
    \item \textbf{Capítulo 4: Metodología.} Aquí se describe el proceso de desarrollo de los modelos y los pasos necesarios para su ejecución.
    \item \textbf{Capítulo 5: Resultados.} Se presentan los hallazgos obtenidos a partir de los modelos desarrollados en el capítulo anterior.
    \item \textbf{Capítulo 6: Discusión y conclusiones.} Aquí se interpretan y analizan los hallazgos de los capítulos previos, relacionándolos con los objetivos planteados y el contexto general.
\end{itemize}
